\documentclass[12pt]{article}
\usepackage[utf8]{inputenc}
\usepackage[T1]{fontenc}
\usepackage[margin=1in]{geometry}
\usepackage{setspace}
\usepackage{xcolor}
\usepackage[colorlinks=true, urlcolor=blue!60!black, linkcolor=blue!60!black]{hyperref}
\setlength{\parindent}{2em}
\setlength{\parskip}{1.25em}
\renewcommand{\baselinestretch}{1.0}
\begin{document}
\begin{center}
{\LARGE \textit{Meyer v. Nebraska}, 262 U.S. 390 (1923)\par}
\vspace{0.8em}
{\large Case Brief --- Four-Page Edition\par}
\end{center}
\vspace{1em}

\section{Memory Jogger}

In a 7--2 decision, the Court struck down Nebraska's post-World War I ban on teaching modern foreign languages to children below the eighth grade, holding that Fourteenth Amendment ``liberty'' protects a teacher's right to teach, a pupil's right to acquire knowledge, and parents' right to direct their children's education (\href{https://www.courtlistener.com/opinion/100233/meyer-v-nebraska/}{CourtListener}).

\section{Detailed Case Facts}

Amid anti-German and nativist sentiment after World War I, Nebraska enacted the Siman Act of 1919. It forbade any person, in any private, denominational, parochial, or public school, to teach any subject in a language other than English, and allowed other languages to be taught as languages only after a pupil passed the eighth grade. Violations were misdemeanors punishable by a fine of \$25 to \$100 or up to thirty days in jail. The Nebraska Supreme Court had read the act to permit ``ancient or dead languages'' such as Latin, Greek, and Hebrew while banning modern languages like German and French.

On May 25, 1920, Robert T. Meyer, an instructor at Zion Parochial School, a Lutheran school in Hamilton County, taught reading in German to ten-year-old Raymond Parpart using Biblical stories. The county attorney entered the classroom and charged him.

\section{Procedural History}

Meyer was convicted in the Hamilton County District Court and fined \$25. The Nebraska Supreme Court affirmed, upholding the act as a police-power measure ensuring ``that the English language should be and become the mother tongue of all children reared in this state.'' On writ of error, the U.S. Supreme Court heard argument on February 23, 1923, and decided the case on June 4, 1923, with the companion cases in \textit{Bartels v. Iowa}.

\section{Judicial Votes}

Reversed, \textbf{7--2}.
\begin{itemize}
\item \textbf{Majority:} Justice McReynolds, joined by Chief Justice Taft and Justices McKenna, Van Devanter, Brandeis, Butler, and Sanford.
\item \textbf{Concurrences:} None.
\item \textbf{Dissent:} Justice Holmes, joined by Justice Sutherland, published in \textit{Bartels v. Iowa}.
\end{itemize}

\section{Holding}

A statute forbidding, under criminal penalty, the teaching of any modern language other than English to children who have not passed the eighth grade arbitrarily infringes the liberty protected by the Due Process Clause and exceeds the police power. The narrow holding rests on the teacher's liberty to teach and the parents' right to engage him. McReynolds's broad catalog of liberties was dicta in 1923, as were his reservations of state power to compel attendance, require English instruction, and prescribe public-school curricula---though the catalog later became foundational.

\section{Analysis of Opinions}

\textbf{Majority (McReynolds).} Citing \textit{Allgeyer}, \textit{Lochner}, and \textit{Adkins}, McReynolds expanded ``liberty'' beyond bodily restraint and contract to acquiring knowledge, marrying, raising children, and worship. Teaching is a ``useful and honorable'' calling, parents have a natural duty and right to control their children's education, and ``[m]ere knowledge of the German language cannot reasonably be regarded as harmful.'' The legislature's wish ``to foster a homogeneous people with American ideals'' was understandable, but the means were forbidden. Invoking Plato's ideal commonwealth and Sparta, where the State reared children, McReynolds called such submergence of the individual alien to American institutions. Because only modern languages were targeted, and in a time of ``peace and domestic tranquility'' no emergency existed, the prohibition was arbitrary.

\textbf{Holmes (dissenting).} With ``hesitation and unwillingness,'' Holmes agreed that a common tongue is a lawful end. Where children hear only Polish, French, or German at home, a legislature could reasonably require that they hear only English at school. The question was one ``upon which men reasonably might differ,'' so the law was not an ``arbitrary fiat.'' Holmes nonetheless joined in striking down Ohio's law because it singled out German.

\section{Examples: Future Applications}

\begin{enumerate}
\item \textbf{Same-Side (Invalidated): Mandatory Public Schooling.} A law requiring all children to attend public schools, outlawing private and parochial schools, violates parental liberty, as \textit{Pierce v. Society of Sisters} (1925) held.
\item \textbf{Same-Side (Invalidated): Heritage-Language Programs.} A ban on private academies teaching Mandarin or Spanish literacy to young children, to promote assimilation, is arbitrary.
\item \textbf{Opposite-Side (Upheld): English Instruction.} A rule that every school teach English reading and writing daily is valid; \textit{Meyer} did not question requiring ``that they shall give instructions in English.''
\item \textbf{Opposite-Side (Upheld): Public Curriculum.} A state may mandate civics and decline to fund foreign-language classes, given ``the State's power to prescribe a curriculum for institutions which it supports.''
\item \textbf{Fence-Sitter: Homeschool Oversight and Opt-Outs.} Annual testing for homeschoolers, or refusing opt-outs from a public-school health curriculum, pits parental liberty against the State's preserved power to regulate education.
\end{enumerate}

\section{Critique}

\textbf{Originalist views.} Critics treat \textit{Meyer} as a \textit{Lochner}-era product that read parental control into ``liberty'' without grounding it in the Clause's original procedural meaning and paved the way for \textit{Griswold}, \textit{Roe}, \textit{Lawrence}, and \textit{Obergefell}. Some libertarian originalists defend the result but would ground it in the Privileges or Immunities Clause or the Ninth Amendment.

\textbf{Progressive views.} Civil-liberties scholars celebrate \textit{Meyer} as a defense of cultural pluralism against xenophobia, later rehabilitated in \textit{Griswold}. Barbara Bennett Woodhouse argues instead that it constitutionalized a patriarchal view of ``the child as property,'' sidelining children's own interests.

\section{Key Quotations}

\begin{quote}
``Without doubt, it denotes not merely freedom from bodily restraint, but also the right of the individual to contract, to engage in any of the common occupations of life, to acquire useful knowledge, to marry, establish a home and bring up children, to worship God according to the dictates of his own conscience, and generally to enjoy those privileges long recognized at common law as essential to the orderly pursuit of happiness by free men.''
\end{quote}
\noindent --- Justice McReynolds, Opinion of the Court, 262 U.S. at 399.

\begin{quote}
``The established doctrine is that this liberty may not be interfered with, under the guise of protecting the public interest, by legislative action which is arbitrary or without reasonable relation to some purpose within the competency of the State to effect.''
\end{quote}
\noindent --- Justice McReynolds, Opinion of the Court, 262 U.S. at 399--400.

\begin{quote}
``That the State may do much, go very far, indeed, in order to improve the quality of its citizens, physically, mentally and morally, is clear; but the individual has certain fundamental rights which must be respected. The protection of the Constitution extends to all, to those who speak other languages as well as to those born with English on the tongue. Perhaps it would be highly advantageous if all had ready understanding of our ordinary speech, but this cannot be coerced by methods which conflict with the Constitution --- a desirable end cannot be promoted by prohibited means.''
\end{quote}
\noindent --- Justice McReynolds, Opinion of the Court, 262 U.S. at 401.

\end{document}
